\title{Melt-pool geometry and solidification time in IN625 laser melting: sensitivity to high-temperature property extrapolation}
\begin{abstract}
Melt-pool dimensions provide accessible calibration targets for thermal models of laser additive manufacturing, yet the solidification history inferred from a fitted field also depends on its thermophysical description. For IN625, the property tables used in the present model stop below the melting interval, leaving the resulting rear temperature profile dependent on extrapolation. Here, an enthalpy-based conduction model of bare-plate laser tracks resolves that dependence through paired rear solidus and liquidus positions. Conductivity and sensible specific heat are extrapolated using either the final tabulated slope or a constant endpoint value, separately and together, while retaining the source-power fraction fitted to one condition's thermographic mean length. The length fit leaves the fusion envelope wider and shallower than the metallographic means. Constant endpoint extrapolation of conductivity lengthens the pool, while the corresponding specific-heat change shortens it. The dominant length increase displaces both rear phase boundaries downstream with much less change in their separation. Across two scan speeds, almost equal separations nevertheless correspond to a passage time through the adopted freezing interval shorter by about one-third at the higher speed. These results connect prescribed property functions to distinct spatial and temporal features of the calibrated field. They identify rear phase-boundary position, separation and material passage time as specific outputs for thermal assessment alongside fusion geometry.
\end{abstract}
